%! Justifying a Claim Based on a Confidence Interval for a Difference of Population Proportions — Unit 6, Lesson 6.9 — Student Packet Cover Sheet
\documentclass[10pt]{article}
\usepackage{apstats-article}
\usepackage{apstats-boxes}
\usepackage{ltablex}
\keepXColumns

\begin{document}

% Full-bleed navy banner
\begin{tikzpicture}[remember picture, overlay]
  \fill[navy] (current page.north west)
    rectangle ([yshift=-0.9in]current page.north east);
\end{tikzpicture}
\vspace{-0.8in}

\begin{center}
  {\color{white}\textbf{\LARGE AP Statistics}}\\[4pt]
  {\color{skymid}\large\textbf{Unit 6}}\\[6pt]
  {\color{white}\textbf{\large Lesson 6.9 \quad Justifying a Claim Based on a Confidence Interval for a Difference of Population Proportions}}
\end{center}

\vspace{0.22in}
\namedateperiod
\vspace{0.18in}

\begin{learningtargetbox}
\small
\begin{enumerate}[leftmargin=1.5em, itemsep=5pt]
  \item I can interpret a two-sample $z$-interval for $p_1 - p_2$ using both population descriptions and the confidence level. \textit{(UNC-4.M, Skill 3.D)}
  \item I can use a two-sample CI to evaluate whether there is convincing evidence of a difference between $p_1$ and $p_2$, including the special role of 0. \textit{(UNC-4.N, Skill 4.B)}
\end{enumerate}
\end{learningtargetbox}

\vspace{0.14in}

\begin{tocbox}
\small
\renewcommand{\arraystretch}{1.5}
\begin{tabularx}{\linewidth}{@{} c l X r @{}}
\rowcolor{sky}
\textbf{\#} & \textbf{Component} & \textbf{Description} & \textbf{Score} \\
\midrule
1 & Warm-Up        & Build a two-sample CI; bridge to interpretation & \blank{1.2cm} \\
2 & Guided Notes   & Interpretation template; the zero rule; worked examples & \blank{1.2cm} \\
3 & Group Activity & Interpret given CIs; evaluate claims about two proportions & \blank{1.2cm} \\
4 & Exit Ticket    & Interpret a CI and justify a claim about $p_1-p_2$ & \blank{1.2cm} \\
5 & Homework       & Interpret and evaluate claims from two-sample CIs & \blank{1.2cm} \\
\midrule
  & \multicolumn{2}{r}{\textbf{Total}} & \blank{1.2cm} \\
\end{tabularx}
\end{tocbox}

\end{document}
